% =============================================================================
% 8 Related Work. Written 2026-10-20 from the approved outline. DAK is cited as
% CONCURRENT work (Ben, 2026-10-20); there is no experimental comparison. The
% OSDI '27 CFP has no concurrent-work window ("Prior or concurrent publication
% in non-peer-reviewed contexts, like arXiv.org ... is permitted"), so the
% explicit discussion below is what supports the label.
% =============================================================================

\section{Related Work}\label{sec:related}

\paragraph{Swapping KV through HBM.} FlexGen places weights, KV and
activations across GPU, CPU and disk for offline, latency-insensitive
batches~\cite{flexgen}. Pie relocates a fixed number of every request's layers
to host memory and prefetches them through HBM buffers ahead of use,
swapping each back out after use; an adaptive expansion factor sets how many
layers move~\cite{pie}. \sysname{} uses the same per-layer prefetch but changes
the unit of placement from the layer to the request. No resident pays the link,
the window is a cache rather than a swap area, and only one token per layer
returns to the host. The axes are orthogonal: Pie's layer split could run
inside our resident tier.
% src: flexgen abstract (quoted in s1). pie s3.1 ("relocate m layers"; "per-
%      layer KV cache as the basic management unit"), s3 ("prefetching data for
%      upcoming layers through a FIFO queue"), s3.3 (swap in and swap out),
%      s3.4 adaptive expansion.

\paragraph{Reading host memory from the kernel.} DirectKV keeps all KV in
CPU memory and has the attention kernel read it directly, tiling through
shared memory to approach HBM-resident performance without staging
buffers~\cite{directkv}. It changes the execution path, not the policy: by its
own account it ``does not change scheduling, batching, or KV-management
policies''~\cite{directkv}. A zero-copy kernel could serve as the rider path
behind \sysname{}'s tier interface (\S\ref{sec:fused}).
% src: directkv abstract ("enabling GPU kernels to directly access CPU-resident
%      KV cache"); s1 ("By tiling through SMEM, we shift bandwidth pressure from
%      the CPU-GPU interconnect to HBM"); s8.1 (quoted).

\paragraph{Concurrent work.} Concurrent with our work, DAK gives the GPU
direct access to host memory through Hopper's Tensor Memory
Accelerator~\cite{dak}. Its attention kernel partitions the KV cache along
the batch dimension, keeping a subset of requests in HBM and the rest in host
memory, so both tiers serve one batch. The split is an input: DAK derives a
global offload ratio from what does not fit in HBM and distributes it across
operations, and it is evaluated on offline batched inference~\cite{dak}.
BOOST places a fixed fraction of every request's KV heads in host memory, in
proportion to the two tiers' bandwidths, so that each wave of thread blocks
reads both tiers at once~\cite{boost}; it rejects placing whole requests on
the host tier because long requests could land there and raise tail
latency~\cite{boost}. \sysname{} makes the split a runtime decision instead.
Every step, a capacity model fed by a measured step time chooses how many and
which whole requests ride, and admission prefers long requests because the
budget is in bytes. Residents keep the engine's stock attention kernel.
% src: dak s5: "It partitions the KV cache along the batch dimension--storing
%      the cache for a subset of requests in local GPU memory and the remainder
%      in remote host memory"; s3: "Given this global offload ratio, DAK will
%      calculate the optimal offloading ratio for each individual operation";
%      s6: "Our experiments are conducted under an offline, batched-inference
%      setup". arXiv v1 28 Apr 2026.
%      boost abstract: "fully utilizing both host and HBM bandwidth requires
%      each wave of GPU threadblocks to access both tiers concurrently and in
%      proportion to their bandwidth ratio"; s4: "Placing an alpha fraction of
%      whole requests on the host tier would satisfy proportionality, but
%      unknown context lengths could let long requests land on the host tier
%      and raise tail latency, so we partition at a finer granularity";
%      DSR "maintains single-head-host and all-local blocks". arXiv 2609.13592.

\paragraph{Moving attention off the GPU.} NEO offloads part of attention and
its KV to the host CPU and balances the two with asymmetric
pipelining~\cite{neo}; its gain shrinks as the GPU grows stronger relative to
the CPU. Model--attention disaggregation runs attention on separate
memory-optimised devices~\cite{lamina}. \sysname{} keeps all compute on the
GPU and uses the host only as memory.
% src: neo abstract ("offloads part of attention compute and KV cache states
%      from the GPU to the local host CPU"; "up to 7.5x, 26%, and 14% ... on
%      T4, A10G, and H100"). lamina abstract v2 ("a collection of cheap,
%      memory-optimized devices for the attention operator").
% TODO-verify: Lamina details beyond the abstract are not used until the PDF
%      is read.

\paragraph{Fetching less.} InfiniGen prefetches only the KV entries it
predicts the next layer will need~\cite{infinigen}, and ShadowKV keeps a
low-rank key cache on the GPU, offloads values, and reconstructs a small
sparse set of KV pairs on the fly~\cite{shadowkv}. Both approximate attention
over a subset of the context; \sysname{} moves every byte and attention stays
exact. The two compose: fewer bytes per rider would let the link carry more
riders.
% src: infinigen abstract: prefetching "only the essential KV cache entries
%      (without fetching them all)". shadowkv abstract: "stores the low-rank key
%      cache and offloads the value cache"; "reconstructs minimal sparse KV
%      pairs on-the-fly".
% TODO-verify: KV compression / quantisation (e.g. KIVI) is not cited until its
%      source is fetched.

\paragraph{Caching KV for reuse.} LMCache moves KV out of the engine so other
engines and later queries can reuse it, across GPU, CPU, storage and
network~\cite{lmcache}. Mooncake builds a disaggregated KV cache from the
cluster's CPU, DRAM and SSD~\cite{mooncake}, and CachedAttention keeps the KV
of past conversation turns to avoid recomputing it~\cite{cachedattention}.
These systems store KV for the next prefill. \sysname{} serves active decode
from host memory; the KV it receives on a decode node could come from such a
store.
% src: lmcache abstract: "shares them across engines and queries"; "across
%      GPU, CPU, storage, and network layers". mooncake abstract: "leverages
%      the underutilized CPU, DRAM, and SSD resources of the GPU cluster";
%      "disaggregated cache of KVCache". cachedattention abstract: "enables
%      reuse of KV caches across multi-turn conversations".
